\chapter{Asking London Billionaires How They Got Rich}

This School of Hard Knocks interview, curated by LazyingArt LLC, is staged as a street search through London, but its inner structure is tighter than the form first suggests. We are led from access to mechanism, from mechanism to a handful of recurring commercial ideas, and from those ideas to a final practical demand for action. The opening montage gives us the endpoints before the reasons: a company sold, a fortune counted in the hundreds of millions, a \$30 million year, a nine-figure company-scale claim, and a \$55 million overnight collapse. The lecture then restarts, and from that point onward it unfolds in a deliberate order: exits first, then agency, then time, then leverage, and finally recurring revenue.

\section{Teaser, field, and access}

The lecture does not begin by teaching. It begins by showing consequences. We hear, in rapid succession, of a company sold for more money than its founder will ever need, of three or four hundred million in wealth, of a company taken to \$150 million and still scaling, of a \$30 million year, and of a \$55 million overnight loss. The first discipline of the notes is not to merge these figures. They belong to different speakers and to different categories.

\begin{definition}
To keep the lecture's categories separate, we shall write
\[
N = \text{net worth}, \qquad
Y = \text{annual earnings}, \qquad
R = \text{revenue},
\]
\[
S = \text{claimed company scale when the category is unclear}, \qquad
T = \text{turnover}, \qquad
V_{\mathrm{exit}} = \text{sale value},
\]
\[
D = \text{deposit}, \qquad
L = \text{loaned amount generated from it}, \qquad
E_q = \text{owner equity}, \qquad
B = \text{borrowing}.
\]
\end{definition}

Only after this teaser does the host reset the scene. London is chosen as a search field because, in the lecture's own reported terms,
\begin{align}
\operatorname{rank}_{\text{wealth}}(\text{London}) &= 6,\\
M_{\text{London}} &> 210{,}000,
\end{align}
where \(M_{\text{London}}\) is the claimed number of millionaires living in the city. The city is therefore not a mere backdrop. It is part of the lecture's argument: if one wants to ask how wealth is made, London is presented as a place where such answers should exist in concentrated form.

Then comes a beat that matters more than it may seem to at first glance. The host asks question after question and gets almost nothing. People turn away, refuse, or keep moving. A passerby finally explains the local style: the wealthy are here, but they are low-key. That refusal sequence is structurally important. It tells us that access is scarce, and it gives the first real answer more weight when it finally comes.

\section{Exit as the first mechanism}

The first businessman gives the lecture its earliest hard claim. He proceeds by negation. One does not get rich by working for someone else. One does not get rich by being paid a salary. One does not get rich merely by starting a business and living inside cash-flow crises. The positive statement is reserved for the end of the sequence: one gets rich by selling the business.

That claim is then pressed. The host does not let it remain a slogan. He says, in effect, that the answer is true but not yet sufficient; what matters is what entrepreneurs get wrong about the exit. Only then does the businessman turn the slogan into a design rule:
\begin{equation}
V_{\mathrm{exit}} \approx \mu T,
\end{equation}
where \(T\) is the turnover a buyer is willing to pay on and \(\mu\) is the multiple implied by the sale. This equation is a note-level reconstruction, not a line written in the lecture, but it captures the commercial mechanism exactly as it is described.

The second point is temporal. The sale is not to be imagined at the end only. It is to be imagined at the beginning:
\begin{equation}
t_{\mathrm{exit}} \approx 4\ \text{years}.
\end{equation}
Again, the four-year horizon is the speaker's heuristic, not a law. What matters is the direction of thought: begin with the likely buyer already in mind.

\subsection{Question \& Answer}

\paragraph{Question.} If starting a business is not itself the payoff, what actually makes one rich?

\paragraph{Answer.} In the lecture's first mechanism, the business is built toward sale. Wealth appears when somebody else wants the firm badly enough to pay for it.

The logic can be written as a short sequence:
\begin{enumerate}
\item identify the class of buyers at the outset;
\item infer what such buyers would value;
\item build the firm's activity so that its turnover \(T\) is intelligible to that buyer;
\item realize wealth when the buyer pays a multiple \(\mu\) on that turnover.
\end{enumerate}

This is the first real tightening of the lecture. The host immediately recaps the point for the audience: do not merely start; start toward a buyer. That recap is important, because once the exit rule has been made explicit, the lecture is free to widen without losing its commercial spine.

\section{Dream, agency, and the hard route}

At this point the lecture could have stayed with M\&A logic. Instead it shifts register through Simon Squibb. The host changes the question from ``How did you get rich?'' to ``What is your dream?'' That is not a detour. It is the lecture's way of asking what sort of human posture is required before any mechanism can work.

Simon's dream, as stated in the interview, is to fix the education system. Young people, he says, are being let down. AI is arriving, old promises about jobs are being repeated too casually, and people are not being given the knowledge that would let them understand they have agency over their own lives. The host then sharpens this into a more immediate social question: if the poor get poorer and the rich get richer, what keeps people broke?

\subsection{Question \& Answer}

\paragraph{Question.} What keeps people broke in today's world?

\paragraph{Answer.} The first answer given here is not leverage, valuation, or even opportunity. It is blame. People blame the rich, blame the government, blame circumstances. Simon's counterclaim is that the first practical move is to blame oneself in the constructive sense: skill up, take agency, and act.

This is where the autobiographical line matters. He says he tried the easy route---benefits, then a job---and that it was the hard route that finally taught him he could create wealth instead of waiting for it. The lecture uses that move very carefully. It is not yet teaching business structure. It is clearing the ground on which business structure will later make sense.

By the end of this segment we are meant to feel a pressure building. If agency is real, then what exactly does the agent need to own? The next section answers that question more economically.

\section{Time, debt, and the ownership of life}

The lecture now asks what the rich do differently. Simon's answer is initially surprising. He does not begin with taxation, product, or market structure. He begins with attitude toward life. The rich, as he describes them, do not think first about retirement, and they do not divide existence into a hated weekday and a liberated weekend. There is no work-life balance in that picture because the point is to build a life in which work is already contained.

The commercial claim becomes sharper once salary enters the frame. Salary can turn into a trap. It does so not because salary is evil in itself, but because debts gather around it: mortgage, car payments, lifestyle commitments. ``Golden handcuffs'' appears in the interview, and then something stronger: a middle-class trap. The lecture's condensed definition is one of its best:
\textit{money buys time.}

That sentence should be read literally. Time is the scarce quantity. To own one's time is to stop pricing one's life only in hours sold to someone else. We may summarize the distinction schematically as
\begin{align}
\text{hours sold} &\longrightarrow \text{income tied directly to time},\\
\text{outcomes sold} &\longrightarrow \text{income tied to result}.
\end{align}
This is not blackboard mathematics from the lecture. It is a careful compression of Simon's claim that he has sold outcomes and results, and has taken a share of success, rather than selling the sheer duration of his labor.

\subsection{Question \& Answer}

\paragraph{Question.} Why, in the lecture's view, is this kind of money knowledge not taught in school?

\paragraph{Answer.} Because once one understands that time is limited and that there are commercial forms other than hourly sale of labor, the ordinary educational pipeline begins to look far less natural.

The lecture states this polemically. Teachers themselves are paid by the hour, and so may not understand the thing they are supposed to teach. Beneath the rhetoric, however, the structural point is clear enough: if time is scarce, then selling time cannot be the only serious model of value.

A short practical interlude follows. The transcript becomes noisy and partly promotional, but the content-bearing point is simple. One should build through a legal structure that separates the person from the business. In England the relevant object is the limited company; in the United States the lecture references the LLC. The lasting point is not the plug. It is that structure protects, accounts, and limits the spillover of business risk into personal liability.

The host then summarizes the whole segment in a more ordinary entrepreneurial phrase: the rich structure their businesses properly. That summary is weaker than Simon's philosophy, but it prepares the lecture to return to exits with greater maturity.

\section{The return to exits}

The earlier invitation to join for a glass of champagne now pays off, and the lecture circles back to the older businessman. This return matters. The first exit lesson was abrupt and defensive. This one is slower, calmer, and much more procedural.

The speaker is in his eighty-sixth year. He speaks of roller-coaster rides, of regrets, and of the surprise of what a lifetime can contain. Only after that widening of time does the host bring back the earlier commercial question: if people want to build a company in order to sell it, how is that actually done?

The answer is not romantic. Join a firm, learn the trade, then build a smaller version of that trade independently, perhaps with one or two colleagues. Acquire clients. Build turnover. Then sell at a multiple. The lecture now gives an explicit rough scale:
\begin{equation}
T \approx 0.5\text{--}1.0\ \text{million}.
\end{equation}
That quantity is not universal and the currency is not forced by the notes. It is simply the interviewee's example of what a saleable business can look like in practice.

\subsection{Question \& Answer}

\paragraph{Question.} How does one actually build a company to sell?

\paragraph{Answer.} Not by inventing in the void, but by learning a trade inside an existing institution, reproducing it independently, building a client base, and making the business legible to buyers.

The path is laid out in sequence:
\begin{enumerate}
\item join a firm and learn the trade;
\item leave with enough knowledge to build a smaller, sharper version;
\item acquire a client base and build turnover \(T\);
\item sell for a multiple, so that \(V_{\mathrm{exit}} \approx \mu T\).
\end{enumerate}

The lecture does something else here that is easy to lose if one compresses too hard. It widens from transaction to character. The speaker's final message is not another business tactic but a general maxim: it is not what fortune throws at us that matters most, but how we react. His Latin motto, \emph{quod tango muto}---that which I touch, I change---recasts entrepreneurship as intervention in the world rather than mere extraction from it.

From here the lecture changes mechanism once again. Having treated exits twice, it moves from sale value to capital structure.

\section{Banking, leverage, and compounding}

The hedge-fund owner shifts the lecture from the economics of selling firms to the economics of financing assets. The segment begins biographically enough: Nigeria, London, investment banking, then a hedge fund of his own. The first hard number is
\begin{equation}
Y_{\text{banker}} \approx \$30 \times 10^6,
\end{equation}
presented as the most money he made in a single year. The category matters: the notes keep this as reported annual earnings, not as net worth and not as fund size.

Then comes an important moral beat before the mathematics. When someone says no, he says, what they often mean is try harder. That line belongs here because the lecture is still preserving its rhythm of resistance followed by penetration. Only after that does the host ask what banks do not want people to know.

\subsection{Question \& Answer}

\paragraph{Question.} What do banks not want people to know?

\paragraph{Answer.} In the lecture's simplified form, a deposit is not passive. It becomes the basis for repeated lending, and therefore for amplified purchasing power.

The spoken heuristic can be written as
\begin{align}
L &\approx 9D,\\
D + L &\approx 10D,
\end{align}
where \(D\) is the initial deposit and \(L\) is the onward lending the banker is describing.

\begin{remark}
This is a transcript-backed rule of thumb, not a precise contemporary banking identity. The lecture uses it to expose the idea of amplification, not to teach regulation.
\end{remark}

Once amplification is in view, leverage follows naturally. The banker says, in effect, that all the visible wealth around us is likely financed: property, restaurants, even the passing car. A minimal notation makes that claim precise:
\begin{align}
A &= E_q + B,\\
\lambda &= \frac{A}{E_q},
\end{align}
where \(A\) is the asset position, \(E_q\) owner equity, \(B\) borrowing, and \(\lambda\) the leverage ratio.

This is the real conceptual pivot of the section. Surface wealth is re-read as structured wealth. Something can be genuinely valuable and still be heavily borrowed against. The lecture wants us to stop confusing possession with unlevered ownership.

The banker ends on a longer horizon. ``Compounding helps you,'' he says. In standard mathematical form, the underlying intuition is
\begin{equation}
W_t = W_0(1+r)^t,
\end{equation}
with the explicit caveat that this equation is our reconstruction of the compounding idea rather than a formula spoken symbolically in the interview. Still, the fit is exact enough. Time magnifies early positions, and so the lecture advises entering rather than obsessing over whether one has paid ten percent too much at the front end.

The host then recaps the whole section in the bluntest possible way: the richest people leverage money. That summary prepares the last speaker, who will connect leverage to something even more operationally important, namely recurring revenue.

\section{Shock, execution, and recurring revenue}

The final major interview begins with velocity. The founder says he made his first million at twenty-three. He is now in health tech, and he says he has brought the company to roughly
\begin{equation}
S \approx \$150 \times 10^6,
\end{equation}
which we preserve as a company-scale claim, not as a forced statement about either revenue or valuation. The ambiguity is part of the source record and should remain visible.

Then comes the lecture's sharpest negative number:
\begin{equation}
\Delta W \approx -\$55 \times 10^6.
\end{equation}
The loss is described as having happened overnight when recession struck, banks foreclosed on facilities, and developments were pulled out from underneath the firm. The important thing is the way the speaker interprets the event. He does not merely call it disastrous. He calls it, in a certain sense, good, because it corrected arrogance.

Before returning to business mechanics, the lecture pauses on something else: execution. Asked what the wealthiest people know that others do not, the founder says everyone knows the formula; the difference is execution. Ideas are cheap. Focus is rare. Sacrifice is real. That beat matters because it prevents the final recurring-revenue section from sounding like a magic trick. It is not a trick. It is a structure that still requires a person who will carry it.

Only then does the host ask the most business-school question in the lecture: how do you scale a \$100 million company?

\subsection{Question \& Answer}

\paragraph{Question.} How does one scale a \$100 million company?

\paragraph{Answer.} In the lecture's closing mechanism, scale comes not from endlessly repeating one-off sales, but from building a revenue stream that recurs, compounds, and can therefore be valued by institutions at a multiple.

Let \(n\) be the number of subscribers and \(p\) the periodic subscription price. Then the note-level reconstruction is
\begin{align}
MRR &= np,\\
ARR &= 12np,\\
V &\approx m \cdot ARR, \qquad m \in [10,20].
\end{align}
Here \(MRR\) is monthly recurring revenue, \(ARR\) is annual recurring revenue, and \(m\) is the valuation multiple the founder says institutions may apply.

The lecture gives one explicit count: \(n=5000\). Keeping the price \(p\) symbolic, we can work through the example exactly as the interview invites us to:
\begin{enumerate}
\item begin with \(n=5000\);
\item then
\[
MRR = 5000p;
\]
\item annualize:
\[
ARR = 12 \cdot 5000p = 60{,}000p;
\]
\item now apply the stated multiple range:
\[
V \approx m \cdot 60{,}000p, \qquad m \in [10,20];
\]
\item hence the implied valuation range is
\[
V \in [600{,}000p,\ 1{,}200{,}000p].
\]
\end{enumerate}

\begin{remark}
The range \(m \in [10,20]\) is a transcript-derived heuristic from the interview, not a universal valuation theorem. The lecture is giving us a business rule of thumb, not a general law of finance.
\end{remark}

This is the chapter's cleanest commercial mechanism. A one-off product sale forces one to go out and sell again. A subscription base keeps sending money while one sleeps. The lecture explicitly adds the second step: institutions then value that recurring stream. A smaller but repeatable revenue base can therefore support a much larger firm value than a more dramatic but non-recurring burst of sales.

The line about the double \(R\) on a Rolls-Royce standing for recurring revenue is best treated as wit rather than doctrine. But the wit lands because the mechanism is now clear. Wealth, in this last segment, is not merely a matter of sales volume. It is a matter of repeatability, compounding, and valuation.

The founder's final exhortation is therefore harsher than motivational boilerplate. Stop stalling. Stop waiting. Stop self-doubting. People less competent than you have made it already. The lecture closes where it wanted to close all along: the structure matters, but structure without action remains inert.

\section{Summary}

The lecture opens with outcomes and only later supplies their mechanisms. First comes London as a wealthy search field and access problem. Then comes the first hard rule: wealth is realized at exit. Simon Squibb widens the frame by arguing that agency comes before technique, and then sharpens it again by identifying money with owned time rather than sold hours. The older businessman returns the lecture to exits, now in a calmer procedural form. The banker turns visible wealth into leverage and compounding. The founder finally joins resilience, execution, and recurring revenue into the lecture's strongest quasi-mathematical explanation.

What remains after the street energy is stripped away is a compact set of claims. Build toward an exit. Refuse the salary-debt trap if you can. Understand that much apparent wealth is leveraged. Prefer revenue streams that recur. And remember that the lecture never lets the mathematics float free from temperament: ideas are cheap, structures are available, but the decisive difference is still the willingness to act.